\chapter{Testarea și Validarea Sistemului StyleGenAI}
\label{cap:testare}

Validarea aplicației StyleGenAI a fost realizată printr-o abordare riguroasă, combinând testarea funcțională (verificarea corectitudinii algoritmilor) cu testarea nefuncțională (performanță, scalabilitate) și testarea de acceptanță (Experiența Utilizatorului - UX).

\section{Metodologia de Testare}
Procesul de asigurare a calității (QA) a urmat ciclul V-Model, validând fiecare componentă față de specificațiile definite în Capitolul 2.

\subsection{Mediul de Testare}
Testele s-au desfășurat într-un mediu controlat cu următoarele specificații:
\begin{itemize}
    \item \textbf{Server Backend}: Intel Core i7-11800H, 16GB RAM, NVIDIA RTX 3060 (pentru inferență AI accelerată CUDA).
    \item \textbf{Client Mobil}: Google Pixel 6 (Android 13) și iPhone 13 (iOS 16, simulator).
    \item \textbf{Conexiune}: Rețea Wi-Fi 5GHz, lățime de bandă 300 Mbps.
\end{itemize}

\section{Matricea Cazurilor de Testare (Test Case Matrix)}
S-au definit 20 de scenarii critice. Tabelul \ref{tab:test_cases} prezintă un eșantion reprezentativ al rezultatelor obținute.

\begin{table}[h]
    \centering
    \begin{tabular}{|l|p{5cm}|p{4cm}|l|}
    \hline
    \textbf{ID} & \textbf{Descriere Scenariu} & \textbf{Rezultat Așteptat} & \textbf{Status} \\ \hline
    TC-01 & Încărcare imagine JPG validă (3MB) & Procesare < 3s, Status 200 OK & \textbf{PASS} \\ \hline
    TC-02 & Încărcare fișier PDF invalid & Eroare 400 "Invalid format" & \textbf{PASS} \\ \hline
    TC-03 & Detecție culoare (Tricou Roșu) & Dominantă: (200, 10, 10) & \textbf{PASS} \\ \hline
    TC-04 & Generare ținută (Ocazie: Nuntă) & Sugestie include sacou/cămașă & \textbf{PASS} \\ \hline
    TC-05 & Autentificare cu parolă greșită & Eroare 401 Unauthorized & \textbf{PASS} \\ \hline
    TC-06 & Testare concurență (10 req/s) & Timp răspuns mediu < 500ms & \textbf{PASS} \\ \hline
    \end{tabular}
    \caption{Rezultatele testării funcționale pentru modulele critice.}
    \label{tab:test_cases}
\end{table}

\section{Evaluarea Performanței Algoritmului AI}

\subsection{Clasificarea Stilului Vestimentar}
Modelul CNN (ResNet-50) a fost evaluat pe un set de date de testare de 500 de imagini noi, neutilizate la antrenare.
Matricea de confuzie indică o acuratețe globală de \textbf{92\%}.

\begin{itemize}
    \item \textbf{Precizie (Precision)}: 0.94 pentru clasa "Casual".
    \item \textbf{Rapel (Recall)}: 0.89 pentru clasa "Formal".
\end{itemize}
Erorile minore au apărut la distincția dintre stilul "Smart Casual" și "Business", due to overlapping features (e.g., cămăși simple fără sacou).

\subsection{Latența Sistemului}
Timpul total de răspuns (Total Round Trip Time - RTT) este compus din latența rețelei, procesarea imaginii și inferența modelului.
Măsurătorile efectuate pe 100 de cereri succesive au relevat:

\begin{itemize}
    \item \textbf{Upload Imagine}: 450 ms (medie pentru 2MB).
    \item \textbf{Pre-procesare (Resize + Normalize)}: 50 ms.
    \item \textbf{Inferență AI (GPU)}: 120 ms.
    \item \textbf{Generare Sugestii (Database Query)}: 30 ms.
    \item \textbf{Total}: $\approx 650$ ms.
\end{itemize}
Rezultatul se încadrează sub pragul de acceptabilitate de 1 secundă pentru interacțiuni non-instante, oferind o experiență utilizator fluentă.

\section{Validarea Experienței Utilizatorului (UAT)}
Un grup de 10 utilizatori beta a testat aplicația timp de o săptămână. Feedback-ul a fost colectat prin chestionarul standard SUS (System Usability Scale).

\textbf{Scorul SUS obținut: 82/100} (încadrat la calificativul "Excelent").
Utilizatorii au apreciat în mod special:
\begin{enumerate}
    \item Simplitatea interfeței ("Clean UI").
    \item Relevanța sugestiilor cromatice.
\end{enumerate}
Ca punct de îmbunătățire, s-a menționat necesitatea adăugării unei funcții de "Editare Manuală" a categoriei hainei detectate eronat.

\section{Concluzii Testare}
Sistemul StyleGenAI demonstrează stabilitate și performanță adecvată pentru un produs MVP (Minimum Viable Product). Toate scenariile critice (Happy Path) au trecut testele funcționale, iar metricile de performanță validează alegerea arhitecturii distribuite și a modelelor CNN optimizate.